\section{Reconstruction from 2D Gaussian Fields}
\label{sec:reconstruction}

The optimizer never sees an image.  It compares two point-evaluable functions.  The teacher
is the frozen field $\field_v$ of \cref{eq:field}.  The student is the current 3D Gaussian
set, rendered at arbitrary image points by a sparse point rasterizer.

\subsection{Point rendering}
\label{sec:point-rasterizer}

Given sample locations $\bm{x} \in \Omega_v$, the point rasterizer projects the 3D Gaussians
into view $v$ exactly as the dense renderer does, using the EWA projection
$P_{iv} = J_{iv} \Sigma_i J_{iv}\T + 0.3\,I_2$ including the $0.3$-pixel dilation, sorts by
view depth, and alpha-composites front to back at each requested point,
\begin{equation}
  \student_v(\bm{x}) \;=\; \sum_{i} T_i(\bm{x})\, \alpha_i(\bm{x})\, \bm{c}_i(\bm{d}_{iv}),
  \qquad
  T_i(\bm{x}) = \prod_{i' < i} \big(1 - \alpha_{i'}(\bm{x})\big),
  \label{eq:student}
\end{equation}
with $\alpha_i(\bm{x}) = o_i \exp(-\tfrac12 q_{iv}(\bm{x}))$, where $q_{iv}$ is the
Mahalanobis form of $P_{iv}$ and $\bm{c}_i(\bm{d}_{iv})$ the view-dependent SH color.  The
point path matches the dense CPU reference renderer on color, alpha, depth, global visible
order, five parameter-gradient families, and retained screen-space gradients within sealed
tolerances, on synthetic grids and on a calibrated 4{,}096-point replacement draw.  This
parity lets us reason about sampled training as subsampled dense
training.\evidence{ara/logic/claims.md C18,
benchmarks/results/20260716_point_rasterizer_parity_RESULT.json}
\todo{Full-image calibrated parity, CUDA point parity, and active off-grid coordinate
gradients are not yet established.  Verify them or scope them out explicitly (T-14).}

\subsection{Training objective}
\label{sec:objective}

Let $p$ be the target risk measure on $\Omega_v$, either uniform over pixels or a continuous
area measure, and let $q$ be a proposal distribution.  Each training step draws a fixed
number $A$ of attempts $\bm{x}_1, \dots, \bm{x}_A \sim q$ and minimizes the
importance-corrected squared error with rejected attempts retained in the denominator,
\begin{equation}
  \widehat{L}_v \;=\;
  \frac{1}{A} \sum_{a=1}^{A}
  \mathbf{1}[\bm{x}_a \text{ accepted}]\;
  \frac{p(\bm{x}_a)}{q(\bm{x}_a)}\;
  \norm{\student_v(\bm{x}_a) - \field_v(\bm{x}_a)}_2^2 .
  \label{eq:sampled-loss}
\end{equation}
Keeping the denominator at the fixed attempt count $A$ rather than the accepted count makes
\cref{eq:sampled-loss} an unbiased estimator of the target risk under rejection-mixture
proposals.  We verified the estimator identity by exact enumeration on finite pixel grids,
including invariance to attempt micro-chunking within
$2\times10^{-12}$.\evidence{ara/logic/claims.md C19}

Two evaluation risks are computed exactly and are never replaced by the training loss.
$\Jpix$ is the mean over all pixel centers and $\Jarea$ a fixed four-offset quadrature per
pixel.  Experiment tables report the validation risks $\Ju$ under uniform sampling and
$\Jq$ under the proposal correction, both against the frozen fields of validation
views.\evidence{benchmarks/compact_only_carrier_stage_ablation.py}

\paragraph{Sampling proposals.}
Teacher-shaped Gaussian proposals are the obvious importance choice.  In a sealed three-seed
comparison they produced no material win over uniform sampling, so the pipeline samples
uniformly.\evidence{ara/logic/claims.md C20}

\subsection{Streaming behaviour}
\label{sec:streaming}

Both sides of \cref{eq:sampled-loss} are point queries, so a training step touches the
compact fields of the sampled views at kilobyte scale, the sampled coordinate batch, the
projected student parameters, and the model and optimizer state.  No dense render target, no
image cache, and no all-view RGB tensor exist anywhere in the loop.  The fields are queried
on demand, and the data boundary of \cref{sec:overview} forbids materializing dense frames.
This is the mechanism behind the memory result, which is quantified only by the protocol of
\cref{sec:memory}.

\subsection{Initialization}
\label{sec:init}

The direct path uses the two standard 3DGS initializations unchanged, since initialization
is independent of the supervision signal.  Random initialization draws points uniformly
inside the camera-derived scene bounds.  SfM initialization loads a sparse COLMAP point
cloud when one exists for the scene.  \todo{Measure both initializations on the
confirmatory scenes.  SfM requires a real COLMAP reconstruction per scene (T-7).}  An
image-free placement alternative that scores rays by querying all fields serves as an
additional control in the ablation study (\cref{sec:init-comparison}).

\begin{TodoBlock}
Freeze the direct-path optimizer configuration before the confirmatory runs, including the
sample budget per step, the view schedule, learning rates, the topology policy, and the
stopping rule.  These are deliberately not chosen in this text (T-3).  Also show a
sampled-versus-dense equivalence curve, that is, quality of compact-supervised training as
a function of samples per step against dense RGB training at matched update counts
(\cref{fig:sample-budget}).
\end{TodoBlock}

\begin{figure}[t]
  \centering
  \figslot{5cm}{%
    \textbf{Figure 4, sample-budget sweep.}  Held-out PSNR on the y-axis versus samples per
    step on a logarithmic x-axis for compact-field supervision, with the dense RGB control
    as a horizontal band at matched update counts.  One panel per scene, seeds as thin
    lines, medians bold.  The purpose is to locate the budget where sampled supervision
    saturates and to justify the chosen default.  Requires the frozen direct-path schedule.}
  \caption{\redcaption{Quality versus sample budget for compact-field supervision.}}
  \label{fig:sample-budget}
\end{figure}
